\section{数据金字塔：真实、仿真与互联网数据}

机器人数据和 LLM 数据最大的不同，是物理接地。LLM 可以从互联网获得海量文本、代码和多模态材料；机器人需要动作、状态、环境和反馈，真实采集昂贵且慢。因此谢晨强调，机器人数据更像金字塔：顶层是真实端侧数据，中层是仿真/合成数据，底层是互联网和人类数据。

\begin{figure}[H]
\centering
\includegraphics[width=0.94\textwidth]{data-pyramid.png}
\caption{机器人数据金字塔：真实数据、仿真数据与互联网/人类数据。}
\end{figure}

\begin{knowledgebox}{读图：三层数据各自解决什么}
顶层真实数据最贴近部署，但数量少、成本高；中层仿真数据可控、可规模化，适合生成长尾场景；底层互联网/人类数据最大规模、最低成本，但物理接地弱。机器人数据策略必须混合三层。
\end{knowledgebox}

\subsection{为什么机器人数据是荒漠}

机器人不像自动驾驶，有上百万辆车每天回传真实道路数据。通用机器人尚未大规模部署，遥操作成本高，任务场景分散，本体差异大。没有端侧规模，就没有像特斯拉那样自然形成的数据引擎。因此机器人需要仿真和合成数据作为前提条件，而不只是锦上添花。

\begin{warningbox}{仿真不是玩具，也不是万能药}
仿真能规模化生成场景和评测，但有 sim-to-real gap。真实数据少、仿真不准、互联网数据弱接地，这三者必须组合，而不是互相替代。
\end{warningbox}

\subsection{本章小结}

机器人数据的核心矛盾是：最真实的数据最少，最多的数据最不真实。数据金字塔是解决这个矛盾的基本框架。

